% R15 component; only input paths and enumerated location/grammar prose are edited.
% Source: revisions/2026-10-04-r15/results/scalar_report/scalar_evidence.tex; commit 1cb3cc9efe135966ec228dfaf51c6841a6dfc97c.
\subsection{Independent scalar training and numerical accuracy}

\label{sec:r15scalarresults}

The separate scalar study uses the original volatility calibration, $(\sigma_I,\sigma_C)=(0.15,0.10)$, and the same one-state capital economy. Eight independent Howard solves generate full-box and matched-tube policies on all four prespecified grids. The policies are then evaluated by actual interpolation-based deployment on 8,192 fresh paths at both 1,024 and 2,048 time steps. The analytical schedule and both frozen scalar neural policies use the same initial states and Brownian paths.

% R15 component; only input paths and enumerated location/grammar prose are edited.
% Source: revisions/2026-10-04-r15/results/scalar_report/table_scalar_training.tex; commit 1cb3cc9efe135966ec228dfaf51c6841a6dfc97c.
\begin{table}[!htbp]
\centering\small
\caption{Independent scalar training and algebraic accuracy}
\label{tab:r15scalartraining}
\begin{tabular}{lrrrrrr}
\toprule
Policy & $N_y/N_t/L$ & $v(-1)$ & $v(0)$ & $v(1)$ & $E_0$ & Seconds \\
\midrule
F1 & 401/256/6 & -3.393994 & -1.314658 & 0.816406 & 1.08E-12 & 3.76 \\
F2 & 801/512/6 & -3.393648 & -1.314316 & 0.817061 & 2.28E-12 & 9.30 \\
F3 & 1601/1024/6 & -3.393474 & -1.314143 & 0.817386 & 5.93E-12 & 24.54 \\
F4 & 2401/1024/9 & -3.393474 & -1.314143 & 0.817386 & 8.94E-12 & 30.26 \\
T1 & 401/256/6 & -3.393994 & -1.314658 & 0.816406 & 1.08E-12 & 3.77 \\
T2 & 801/512/6 & -3.393648 & -1.314316 & 0.817061 & 2.28E-12 & 9.33 \\
T3 & 1601/1024/6 & -3.393474 & -1.314143 & 0.817386 & 5.93E-12 & 24.64 \\
T4 & 2401/1024/9 & -3.393474 & -1.314143 & 0.817386 & 8.94E-12 & 30.50 \\
\bottomrule
\end{tabular}
\par\medskip\noindent\footnotesize F1--F4 solve the full primitive action box; T1--T4 solve the common schedule tube of radius $0.1$. $E_0$ is the outward finite-grid algebraic error bound in Proposition~\ref{prop:r15scalarresidual}. The displayed error endpoints are rounded upward; the full outward endpoints are retained in the artifact. Seconds include each new solve, its interval residual audit, and policy serialization. The last grid enlarges the domain from $[-6,6]$ to $[-9,9]$ at the same declared spatial and temporal spacing. Boundary data are the stated far-state approximation.
\end{table}


All 8 classical solves reached the stated stopping criterion, using at most 3 Howard updates at each date. Their finite-grid algebraic bounds range from $1.072\times10^{-12}$ to $8.938\times10^{-12}$. The last joint state--time refinement changes the three initial-state values by amounts between $1.724\times10^{-4}$ and $3.250\times10^{-4}$. The subsequent enlargement of the domain changes these values by at most $4.441\times10^{-16}$. Thus the algebraic solve, the training-grid approximation, and the finite-domain approximation have separate, observable numerical accounts. The small domain change does not bound the remaining time or spatial discretization error.

% R15 component; only input paths and enumerated location/grammar prose are edited.
% Source: revisions/2026-10-04-r15/results/scalar_report/table_scalar_deployment.tex; commit 1cb3cc9efe135966ec228dfaf51c6841a6dfc97c.
\begin{table}[!htbp]
\centering\small
\caption{Deployed scalar policy gains on common diffusion paths}
\label{tab:r15scalardeployment}
\begin{tabular}{lrrrrr}
\toprule
Policy & $y=-1$ & $y=0$ & $y=1$ & Population & MC s.e. \\
\midrule
Schedule & 0.0000 & 0.0000 & 0.0000 & 0.0000 & 0.0000 \\
F1 & 0.1137 & 1.5586 & 0.6768 & 0.7818 & 0.0069 \\
F2 & 0.1137 & 1.5587 & 0.6771 & 0.7819 & 0.0069 \\
F3 & 0.1136 & 1.5586 & 0.6772 & 0.7819 & 0.0069 \\
F4 & 0.1136 & 1.5586 & 0.6772 & 0.7819 & 0.0069 \\
T1 & 0.1137 & 1.5586 & 0.6768 & 0.7818 & 0.0069 \\
T2 & 0.1137 & 1.5587 & 0.6771 & 0.7819 & 0.0069 \\
T3 & 0.1136 & 1.5586 & 0.6772 & 0.7819 & 0.0069 \\
T4 & 0.1136 & 1.5586 & 0.6772 & 0.7819 & 0.0069 \\
NBO (R12) & 0.1036 & 1.5557 & 0.5683 & 0.7412 & 0.0071 \\
DPO (R12) & 0.0468 & 1.4460 & 0.6303 & 0.7065 & 0.0066 \\
\bottomrule
\end{tabular}
\par\medskip\noindent\footnotesize All payoff gains and standard errors are multiplied by $10^3$. The comparator is the analytical schedule evaluated on the same numerical mesh and Brownian path. Population denotes the mean over independent uniform initial-state draws; the last column is its paired Monte Carlo standard error. Initial-state columns retain their separate rankings. NBO and DPO are the prespecified frozen R12 scalar policies. These are conditional Euler-deployment diagnostics; Monte Carlo standard errors do not include continuous-time or boundary approximation error.
\end{table}


At 2,048 deployment steps, the largest-domain classical policy has a mean schedule gain of $7.819\times10^{-4}$, compared with $7.412\times10^{-4}$ for NBO and $7.065\times10^{-4}$ for DPO. The paired classical advantages are $4.069\times10^{-5}$ (Monte Carlo standard error $6.841\times10^{-7}$) against NBO and $7.541\times10^{-5}$ ($4.063\times10^{-7}$) against DPO. Both advantages are positive at each of the three reported initial states. The numerical comparison therefore identifies an independently computed feasible classical competitor with higher observed payoffs than both frozen neural policies in this scalar economy.

The neural ordering varies across initial states. NBO minus DPO has population mean $3.472\times10^{-5}$, but its conditional mean at $y=1$ is $-6.206\times10^{-5}$. The same reversal occurs on the 1,024-step deployment mesh. This state-specific outcome is retained alongside the population comparison. Corresponding full-box and tube policies have identical payoffs path by path on both reported deployment meshes. The maximum recorded nodal deviation from the schedule is 0.046600, below the tube radius. This observation concerns the computed scalar policies and does not establish that the tube restriction is harmless in every dimension.

The complete full-box and matched-tube value arrays coincide bit for bit on each of the four grids. Because the action sets are nested and the terminal and boundary data are identical, Corollary~\ref{cor:r15scalartube} converts the two algebraic enclosures into a bound on the loss from the tube restriction in the exact finite-grid problem. Across the four grids, this loss at every initial-date node is at most $1.788\times10^{-11}$. This deterministic conclusion concerns the stated finite-grid problem; the domain and diffusion approximation accounts retain their separate scope.

The complete scalar process consumed 195.90 seconds, including fresh-process setup, all eight solves, interval audits, serialization, both common-path deployments, and output. The individual new classical solves consumed 3.76--30.50 seconds. Within the fine-mesh deployment, the measured action-evaluation totals were 1.26--1.36 seconds for the classical tables and 9.38--9.60 seconds for the frozen neural maps. These are operation-specific measurements on the recorded worker. The historical neural fitting clocks are not included in a new end-to-end training comparison.

The algebraic error bound concerns the exact equations with their declared boundary data. The independent domain enlargement and paired deployment-mesh comparisons quantify distinct numerical sensitivities. The complete tables retain all trained policies, both meshes, all initial-state strata, and all 110 pairwise numerical comparisons. They supply a scalar accuracy and classical-work benchmark without identifying an exact value for the high-dimensional diffusion.
